% !TEX root = ../discretemath.tex

\section{Примеры для симметричной версии локальной леммы Ловаса}

\subsection{Напоминание: симметричная форма}

\begin{Theorem}{Симметричная форма ЛЛЛ (формулировка)}{}
	Пусть $\Pr(A_i)\le q$ для всех $i$, и каждое событие зависит не более чем от $d$ других. Если
	\[
	\mathrm e\,q(d+1)\le 1
	\qquad(\text{грубое достаточное условие: } 4qd\le 1),
	\]
	то $\Pr\!\left(\overline{A_1}\wedge\dots\wedge\overline{A_n}\right)>0$ — с положительной вероятностью не происходит ни одно «плохое» событие. (Доказательство — в билете о локальной лемме Ловаса.)
\end{Theorem}

\subsection{Выполнимость $k$-КНФ}

\begin{Example}{$k$-КНФ}{}
	Рассмотрим $k$-КНФ $C_1\wedge C_2\wedge\dots\wedge C_m$, где каждый дизъюнкт $C_i$ содержит ровно $k$ литералов. Выберем значение каждой переменной независимо и равновероятно ($0$ или $1$). Обозначим
	\[
	A_i=\{\text{дизъюнкт $C_i$ не выполнен}\}.
	\]
	Дизъюнкт из $k$ литералов не выполнен лишь при единственном наборе значений своих переменных, поэтому
	\[
	\Pr(A_i)=2^{-k}.
	\]

	События $A_i$ и $A_j$ зависимы только если $C_i$ и $C_j$ имеют общую переменную. Если каждый дизъюнкт делит переменные не более чем с $2^{k-2}$ другими дизъюнктами, возьмём
	\[
	q=2^{-k},\qquad d\le 2^{k-2}.
	\]
	Тогда
	\[
	4qd\le 4\cdot 2^{-k}\cdot 2^{k-2}=1,
	\]
	и по грубой форме ЛЛЛ
	\[
	\Pr\!\left(\overline{A_1}\wedge\dots\wedge\overline{A_m}\right)>0.
	\]
	Значит, существует набор значений переменных, при котором выполнены все дизъюнкты: такая $k$-КНФ выполнима.
\end{Example}


\subsection{Независимое множество из списков}

\begin{Example}{Независимое множество из списков}{}
	Пусть $G$ — граф с максимальной степенью не больше $d$. Даны попарно непересекающиеся множества вершин $V_1,V_2,\dots,V_t$ с $|V_i|\ge 2\mathrm e d$ для всех $i$. Докажем, что можно выбрать по одной вершине $u_i\in V_i$ так, чтобы $\{u_1,\dots,u_t\}$ было независимым.
\end{Example}

\begin{figure}[H]
	\centering
	\begin{tikzpicture}[font=\small, >=Stealth,
		v/.style={circle, draw, fill=black!80, inner sep=0pt, minimum size=4pt},
		sel/.style={circle, draw, fill=red!80, inner sep=0pt, minimum size=5pt},
		grp/.style={rounded corners, draw, fill=black!4}]
		\foreach \i/\x in {1/0, 2/2.4, 3/4.8}{
			\draw[grp] (\x-0.5,-1.3) rectangle (\x+0.5,1.3);
			\node at (\x,1.6) {$V_{\i}$};
			\foreach \y in {-0.9,-0.3,0.3,0.9}{ \node[v] (n\i\y) at (\x,\y) {}; }
		}
		% выбранные вершины (по одной из каждого списка)
		\node[sel] (s1) at (0,0.3) {};
		\node[sel] (s2) at (2.4,-0.3) {};
		\node[sel] (s3) at (4.8,0.9) {};
		\node at (2.4,-1.7) {выбрано по одной $u_i\in V_i$, рёбер между ними нет};
	\end{tikzpicture}
	\caption{Из каждого списка $V_i$ выбирается одна вершина. «Плохое» событие $A_r$ — оба конца ребра $r$ оказались выбранными. ЛЛЛ гарантирует, что при $|V_i|\ge 2\mathrm e d$ с положительной вероятностью ни одно $A_r$ не происходит — выбор независим.}
\end{figure}

\subsubsection*{Применение симметричной ЛЛЛ}

Положим $K=\lfloor 2\mathrm e d\rfloor+1$ и уменьшим каждое $V_i$ до размера $K$. Из каждого $V_i$ независимо и равновероятно выберем одну вершину. Для каждого ребра $r$ графа $G$ с концами в разных списках определим событие
\[
A_r=\{\text{выбраны оба конца ребра $r$}\},
\qquad
\Pr(A_r)=\frac{1}{K^2}.
\]

Событие $A_r$ зависит лишь от событий, связанных с теми же двумя списками, что и концы $r$. Если $r$ соединяет $V_i$ и $V_j$, то число рёбер с хотя бы одним концом в $V_i\cup V_j$ не превосходит $2Kd$; исключая само $r$, число зависимостей
\[
D\le 2(Kd-1).
\]

Проверим условие симметричной ЛЛЛ:
\[
\mathrm e\cdot \frac{1}{K^2}\cdot (D+1)
\le \mathrm e\cdot \frac{2(Kd-1)+1}{K^2}
<\frac{2\mathrm e d}{K}<1,
\]
поскольку $K>2\mathrm e d$. Следовательно, с положительной вероятностью не происходит ни одно $A_r$: выбранные вершины попарно не соединены рёбрами. Значит, существует независимое множество, содержащее ровно по одной вершине из каждого $V_i$.
